\chapter{Weighted Monge--Amp\`ere measures}
\label{chap:damped-measures}

An extended Monge--Amp\`ere measure may have infinite mass. Multiplying it by a weight that vanishes sufficiently fast at $-\infty$ gives finite measures whose weak limits can be controlled by convergence of the potentials in capacity. Recovering the undamped measure requires a further condition: mass must not persist on increasingly deep sublevels. This condition is necessary even for a sequence of bounded potentials on a K\"ahler curve.

For nonpositive $u\in\PSH(X,\omega)$, write
\[
 \mu_u=\mathrm{MA}^{\mathrm{ext}}_\omega(u),
 \qquad r_n=\max(n-2,0).
\]
A Gauduchon trace bound and Stokes' theorem give mass growth of order $(1+A)^{r_n}$ for potentials of range $A$. This estimate makes exponential weights and polynomial weights of order $q>r_n$ integrable, uniformly over the potentials. Truncation then reduces continuity to the bounded-product calculus of \cref{sec:psh-bounded-products,sec:hermitian-capacity}. The arguments first use the positive background $\omega$; domination by a larger Hermitian form extends the measure estimates to a general smooth background.

\section{Mass growth and the range of a potential}

The local cutoff estimate gives a polynomial bound in the range. The compact geometry improves its exponent: a uniform bound for the first mixed mass removes one power, and the exactness of the top Hessian wedge removes another.

\begin{lemma}\label{lem:bounded-range-mass}
There is $C_\omega<+\infty$ such that, for every $A\geq0$ and $v\in\PSH(X,\omega)$ with $0\leq v\leq A$,
\[
 \int_X(\omega+\ddc v)^n\leq C_\omega(1+A)^n.
\]
\end{lemma}
\begin{proof}
Choose finitely many coordinate balls with smaller relatively compact balls covering $X$. On each larger ball choose a fixed quadratic function $q$ with $\ddc q\geq\omega$. Then $p=v+q$ is plurisubharmonic, and, after a constant shift of $q$, satisfies $0\leq p\leq A+C$ there. Let $\kappa$ be the Euclidean K\"ahler form.

By \eqref{eq:local-cln-range}, for $1\leq j\leq n$ the integral of $(\ddc p)^j\wedge\kappa^{n-j}$ on a fixed smaller ball is at most $C_j(1+A)^j$.

Finally expand
\[
 (\omega+\ddc v)^n
 =\sum_{j=0}^n\binom nj(\ddc p)^j\wedge(\omega-\ddc q)^{n-j}.
\]
Each smooth factor has uniformly bounded coefficients on the ball, so its absolute contribution is bounded by a fixed multiple of the corresponding local integral above. Sum over $j$ and the finite cover.
\end{proof}

The next estimate uses the uniform trace bound from \cref{cor:trace-bounds}.

\begin{proposition}\label{prop:gauduchon-range-growth}
For every Hermitian metric $\omega$ there is a constant $C_\omega$ such that
\[
 \int_X(\omega+\ddc v)^n\leq C_\omega(1+A)^{r_n}
 \quad\text{if }v\in\PSH(X,\omega),\ 0\leq v\leq A.
\]
\end{proposition}
\begin{proof}
We retain the fixed coordinate balls, $q$, $p=v+q$, and $\kappa$ from the preceding proof. The first local Hessian integral has a bound independent of $A$. Indeed,
\[
 \ddc p=(\omega+\ddc v)+(\ddc q-\omega),
\]
and the integral of the first positive term against $\kappa^{n-1}$ on any one coordinate ball is bounded by a fixed multiple of
\[
 \int_X(\omega+\ddc v)\wedge\omega^{n-1}.
\]
This is uniformly bounded by \cref{cor:trace-bounds}; the smooth second term has fixed mass. For $n=1$ the same assertion follows directly from Stokes' theorem.

The first inequality in \eqref{eq:local-cln-range} therefore gives, for $1\leq j\leq n$ on a fixed smaller ball,
\[
 \int(\ddc p)^j\wedge\kappa^{n-j}\leq C_j(1+A)^{j-1}.
\]
Expand the positive mixed measures in these local corrected Hessians and sum over the finite cover. With
\[
 M_j(v)=\int_X(\omega+\ddc v)^j\wedge\omega^{n-j},
 \qquad 0\leq j\leq n,
\]
we obtain
\[
 0\leq M_j(v)\leq C_j(1+A)^{\max(j-1,0)}.
\]

There is one more global cancellation. For the bounded quasi-plurisubharmonic function $v$ and $1\leq j\leq n$, the signed products $(\ddc v)^j$ are globally defined by the same local Bedford--Taylor expansion. They are closed and satisfy
\[
 (\ddc v)^j=\ddc\bigl(v(\ddc v)^{j-1}\bigr).
\]
These identities follow locally from the bounded Bedford--Taylor recursion after a smooth correction, and hence agree on overlaps. Stokes' theorem therefore gives
\[
 0=\int_X(\ddc v)^n
   =\sum_{j=0}^n(-1)^{n-j}\binom nj M_j(v).
\]
Since all $M_j(v)$ are nonnegative,
\[
 M_n(v)\leq\sum_{j=0}^{n-1}\binom nj M_j(v)
          \leq C_\omega(1+A)^{\max(n-2,0)}.
\]
This proves the proposition, including dimension one.
\end{proof}

In dimensions one and two, $r_n=0$, so \cref{prop:gauduchon-range-growth} gives bounded mass. In higher dimensions the exponent is a uniform bound; its optimal value may depend on the metric.

\section{Exponential damping}

An exponential transform has bounded range and preserves positivity. Its Hessian dominates an exponential multiple of the original positive current, so its top product controls the weighted measure.

For $a>0$ set
\[
 F_a(t)=a^{-1}\mathrm{e}^{at}\quad(t\leq0).
\]
Then $0<F_a\leq a^{-1}$, $0<F_a'\leq1$, and $F_a''\geq0$. We set $F_a(-\infty)=0$.

\Needspace{7\baselineskip}
\begin{proposition}\label{prop:exponential-transform}
If $u\in\PSH(X,\omega)$ and $u\leq0$, then $F_a(u)$ is a bounded $\omega$-plurisubharmonic function, and
\[
 \mathrm{e}^{na u}\mu_u\leq\mathrm{MA}_\omega(F_a(u)).
\]
\end{proposition}
\begin{proof}
First suppose $u$ is bounded. The chain rule \eqref{eq:bounded-convex-chain-rule} gives
\[
 \omega+\ddc F_a(u)
 =\mathrm{e}^{au}(\omega+\ddc u)+(1-\mathrm{e}^{au})\omega
      +a\mathrm{e}^{au}\,\mathrm{d}u\wedge\dc u.
\]
All three summands are positive. Positivity of their bounded products yields the claimed measure inequality.

For arbitrary $u$, use $u_t=\max(u,-t)$. The bounded potentials $F_a(u_t)$ decrease and lie between zero and $a^{-1}$, so their limit $F_a(u)$ is a bounded $\omega$-plurisubharmonic function. On the plurifine open set $D_t=\{u>-t\}$, it agrees with $F_a(u_t)$. Locality and the bounded case therefore give
\[
 \mathrm{e}^{na u}\mathbf1_{D_t}\mu_u
 \leq\mathbf1_{D_t}\mathrm{MA}_\omega(F_a(u)).
\]
Let $t\to+\infty$. The pole set is discarded by $\mu_u$ and has zero mass for the bounded product on the right.
\end{proof}

\Needspace{12\baselineskip}
\begin{theorem}\label{thm:exponential-damping}
For every $p>0$, the measure
\[
 \nu_{p,u}:=\mathrm{e}^{pu}\mu_u
\]
is finite, non-pluripolar, and satisfies
\[
 \nu_{p,u}(X)\leq C_\omega(1+p^{-1})^{r_n}
 \quad\text{for all }u\in\PSH(X,\omega),\ u\leq0.
\]
The constant depends only on $\omega$.
\end{theorem}
\begin{proof}
Take $a=p/n$ in \cref{prop:exponential-transform} and apply \cref{prop:gauduchon-range-growth} with $A=n/p$. The measure vanishes on pluripolar sets because $\mu_u$ does.
\end{proof}

Since $\mu_u$ vanishes on the pole set, multiplication by $\mathrm{e}^{-pu}$ recovers it from $\nu_{p,u}$, with infinite mass allowed. An additive shift $u\mapsto u+c$ multiplies the weighted measure by $\mathrm{e}^{pc}$.

\section{Capacity bounds for damped measures}

The exponential transform becomes a test for $\operatorname{Cap}_\omega$ after rescaling its range to $[0,1]$. The resulting estimate bounds concentration uniformly over the nonpositive potentials. Fix $p>0$.

\begin{proposition}\label{prop:damped-capacity}
For every Borel set $E$,
\[
 \nu_{p,u}(E)
 \leq\max(1,n/p)^n\operatorname{Cap}_\omega(E).
\]
\end{proposition}
\begin{proof}
Put $A=\max(1,n/p)$ and $v=F_{p/n}(u)$. Then $0\leq v/A\leq1$ and $v/A$ is $\omega$-plurisubharmonic. Moreover
\[
 A(\omega+\ddc(v/A))=(\omega+\ddc v)+(A-1)\omega
 \geq\omega+\ddc v.
\]
Hence $\mathrm{MA}_\omega(v)\leq A^n\mathrm{MA}_\omega(v/A)$. Combine this with \cref{prop:exponential-transform} and the definition of capacity.
\end{proof}

\section{Polynomial damping}

Polynomial damping follows from the same range estimate by integrating over the level parameter. The first step is to bound the mass on a superlevel set where the potential has bounded depth.

\begin{lemma}\label{lem:sublevel-mass-growth}
There is a constant $C$ depending only on $\omega$ such that
\[
 \mu_u(\{u\geq-t\})\leq C(1+t)^{r_n}
 \quad(t\geq0, u\in\PSH(X,\omega),\ u\leq0).
\]
\end{lemma}
\begin{proof}
Put $v_t=\max(u,-t-1)+t+1$. Then $0\leq v_t\leq t+1$ and $v_t$ is $\omega$-plurisubharmonic. On $\{u\geq-t\}\subset\{u>-t-1\}$ its measure agrees with $\mu_u$ by plurifine locality. Thus the required mass is at most $\int_X(\omega+\ddc v_t)^n$, which \cref{prop:gauduchon-range-growth} bounds by $C(1+t)^{r_n}$.
\end{proof}

\Needspace{15\baselineskip}
\begin{theorem}\label{thm:polynomial-damping}
For every $q>r_n$, the measures
\[
 \sigma_{q,u}:=(1-u)^{-q}\mu_u
\]
have uniformly bounded mass for $u\in\PSH(X,\omega)$, $u\leq0$. More precisely,
\[
 \sigma_{q,u}(\{u<-T\})
 \leq C\frac{q}{q-r_n}(1+T)^{r_n-q}\quad(T\geq0).
\]
There is also $C_q$ such that, for every Borel set $E\subset X$,
\[
 \sigma_{q,u}(E)\leq C_q\operatorname{Cap}_\omega(E)^{(q-r_n)/(q+n-r_n)}.
\]
\end{theorem}
\begin{proof}
Write $s=-u\geq0$. The identity
\[
 (1+s)^{-q}=q\int_s^\infty(1+t)^{-q-1}\,\mathrm{d}t
\]
and Tonelli's theorem, valid also for the extended measure $\mu_u$, give the total-mass bound from \cref{lem:sublevel-mass-growth}. Restricting to $\{s>T\}$ gives
\[
 \begin{split}
 \int_{\{s>T\}}(1+s)^{-q}\,\mathrm{d}\mu_u
 &\leq q\int_T^\infty(1+t)^{-q-1}\mu_u(\{s\leq t\})\,\mathrm{d}t\\
 &\leq C\frac{q}{q-r_n}(1+T)^{r_n-q}.
 \end{split}
\]
For the capacity estimate write $\delta=\operatorname{Cap}_\omega(E)$. If $\delta=0$, the set $E$ is pluripolar, so $\sigma_{q,u}(E)=0$. If $0<\delta\leq1$, choose $T=\delta^{-1/(q+n-r_n)}\geq1$. Put $v_T=\max(u,-T-1)+T+1$. Since $v_T/(T+1)$ is an admissible capacity test, positivity and plurifine locality give
\[
 \sigma_{q,u}(E\cap\{s\leq T\})
 \leq\mu_u(E\cap\{s\leq T\})
 \leq (1+T)^n\delta.
\]
The remaining part is bounded by the preceding tail estimate. Since $T^n\delta=T^{r_n-q}=\delta^{(q-r_n)/(q+n-r_n)}$, both terms have the claimed bound. For $\delta>1$, use the uniform total-mass bound.
\end{proof}

For $n\geq2$, the capacity exponent is $1-n/(q+2)$. The threshold $q>n-2$ comes from the mixed-mass estimate in \cref{prop:gauduchon-range-growth}.

\Needspace{14\baselineskip}
\section{Continuity of the damped operators}

At a fixed truncation level, capacity convergence gives weak convergence of the bounded products and their continuous weights. The uniform weighted tail estimates allow the truncation level to tend to infinity. Here and below, weak convergence of finite measures is tested against continuous functions on $X$.

\begin{theorem}\label{thm:damped-continuity}
Let $u_j,u\in\PSH(X,\omega)$ be nonpositive, and suppose $u_j\to u$ in capacity. Then, for every $p>0$ and every $q>r_n$,
\[
 \mathrm{e}^{pu_j}\mu_{u_j}\rightharpoonup \mathrm{e}^{pu}\mu_u,
 \qquad
 (1-u_j)^{-q}\mu_{u_j}\rightharpoonup(1-u)^{-q}\mu_u.
\]
The undamped measures are allowed to have infinite mass.
\end{theorem}
\begin{proof}
For $T>1$, choose a continuous function $\chi_T:\R\to[0,1]$ which is zero on $(-\infty,-2T]$ and one on $[-T,+\infty)$. Let $w(t)$ be either $\mathrm{e}^{pt}$ or $(1-t)^{-q}$ on $t\leq0$. Plurifine locality gives
\[
 \chi_T(u_j)w(u_j)\mu_{u_j}
 =g_T(\max(u_j,-2T))\,
       \mathrm{MA}_\omega(\max(u_j,-2T)),
\]
where $g_T(t)=\chi_T(t)w(t)$ on $[-2T,0]$ and $g_T(-2T)=0$.
The truncated potentials are uniformly bounded and converge in capacity to the corresponding truncation of $u$, since truncation is Lipschitz. Since $g_T$ is uniformly continuous on $[-2T,0]$, the coefficients $g_T(\max(u_j,-2T))$ also converge in capacity and are uniformly bounded. The bounded product measures are uniformly dominated by a fixed multiple of capacity, using the capacity test obtained by shifting and scaling the truncations. The integral of the difference of the coefficients therefore tends to zero: outside a set of capacity tending to zero the difference is uniformly small. For the limiting quasi-continuous coefficient, remove a set of arbitrarily small capacity and extend its continuous restriction from the compact complement. Bounded Bedford--Taylor continuity applies to this continuous test function, and capacity domination bounds the removed part. This proves weak convergence of the truncated weighted measures for each fixed $T$.

It remains to remove the cutoff uniformly in $j$. For exponential damping,
\[
 \int_{\{u_j<-T\}}\mathrm{e}^{pu_j}\,\mathrm{d}\mu_{u_j}
 \leq \mathrm{e}^{-pT/2}\int_X\mathrm{e}^{pu_j/2}\,\mathrm{d}\mu_{u_j}
 \leq C_p \mathrm{e}^{-pT/2}
\]
by \cref{thm:exponential-damping}. For polynomial damping, use \cref{thm:polynomial-damping}. Both bounds tend to zero and apply also to $u$. Testing against a bounded continuous function, pass first $j\to\infty$ at fixed $T$ and then $T\to\infty$. This proves both assertions.
\end{proof}

Choose the potential $U$ used in \cref{cor:strong-linear-infinite-mass}. For $0<\varepsilon<1$, the potentials $\varepsilon U$ approach zero in capacity as $\varepsilon\downarrow0$. This follows from quasi-continuity and the capacity-zero pole set: outside a set of arbitrarily small capacity, $U$ is finite, continuous and bounded on a compact set. Hence
\[
 \mathrm{e}^{p\varepsilon U}\mu_{\varepsilon U}\rightharpoonup\omega^n,
 \qquad
 (1-\varepsilon U)^{-q}\mu_{\varepsilon U}\rightharpoonup\omega^n,
\]
although $\mu_{\varepsilon U}(X)=+\infty$ for every $0<\varepsilon<1$. The full currents also converge to $\omega$ in the coefficient $L^1$ norm, and their Gauduchon traces converge in total variation.

\Needspace{17\baselineskip}
\begin{theorem}\label{thm:damped-continuity}\label{prop:general-damping-weights}
Let $w:[0,+\infty)\to(0,+\infty)$ be continuously differentiable and nonincreasing, with $w(t)\to0$. For $T\geq0$ define
\[
 B_w(T)=\int_T^\infty(1+t)^{r_n}(-w'(t))\,\mathrm{d}t.
\]
Assume $B_w(0)<+\infty$. For nonpositive $u\in\PSH(X,\omega)$, the measures
$\nu_{w,u}=w(-u)\mu_u$, with weight zero on the pole set, satisfy
\[
 \nu_{w,u}(X)\leq C_\omega B_w(0),\qquad
 \nu_{w,u}(\{u<-T\})\leq C_\omega B_w(T).
\]
For every Borel set $E$ and $T\geq0$,
\[
 \nu_{w,u}(E)\leq
 w(0)(1+T)^n\operatorname{Cap}_\omega(E)+C_\omega B_w(T).
\]
In particular they are uniformly absolutely continuous with respect to capacity. If nonpositive $u_j,u\in\PSH(X,\omega)$ satisfy $u_j\to u$ in capacity, then $\nu_{w,u_j}\rightharpoonup\nu_{w,u}$.
\end{theorem}
\begin{proof}
Write $s=-u$. Since $w(s)=\int_s^\infty(-w'(t))\,\mathrm{d}t$, Tonelli's theorem and \cref{lem:sublevel-mass-growth} give
\[
 \int_X w(s)\,\mathrm{d}\mu_u
 =\int_0^\infty(-w'(t))\mu_u(\{s\leq t\})\,\mathrm{d}t
 \leq C_\omega B_w(0).
\]
This identity includes the set $\{s=0\}$; the pole set has zero $\mu_u$-mass. Restricting the integral to $\{s>T\}$ gives the same bound with $B_w(T)$, since the inner set is then $\{T<s\leq t\}$ and $t\geq T$.

The second bound in \eqref{eq:restricted-mass-growth} gives
\[
 \int_{E\cap\{s\leq T\}}w(s)\,\mathrm{d}\mu_u
 \leq w(0)(1+T)^n\operatorname{Cap}_\omega(E).
\]
Add the tail estimate. Since $B_w(T)\to0$, first taking $T$ large and then the capacity small gives uniform absolute continuity.

For continuity, fix a continuous nonincreasing $\eta:[0,+\infty)\to[0,1]$, equal to one on $[0,1]$ and zero on $[2,+\infty)$, and put $\eta_T(s)=\eta(s/T)$. For fixed $T>0$, set $v_{j,T}=\max(u_j,-2T)$ and $g_T(t)=\eta_T(-t)w(-t)$ on $[-2T,0]$. Since $g_T(-2T)=0$, locality gives
\begin{equation}\label{eq:weighted-truncation-identity}
 \eta_T(-u_j)w(-u_j)\mu_{u_j}
 =g_T(v_{j,T})\,\mathrm{MA}_\omega(v_{j,T}).
\end{equation}
Truncation preserves capacity convergence. The functions $g_T(v_{j,T})$ are uniformly bounded and quasi-continuous, and converge in capacity to the corresponding coefficient for $u$. The weighted bounded-product continuity in \cref{sec:psh-local-capacity,sec:hermitian-capacity} therefore gives weak convergence in \eqref{eq:weighted-truncation-identity} at each fixed $T$. The tail bounds control the error by $C_\omega B_w(T)$ times the supremum norm of a continuous test function, uniformly in $j$ and for $u$. Let $T\to\infty$ to conclude.
\end{proof}

For $w(t)=\mathrm{e}^{-pt}$, $p>0$, the integral $B_w(0)$ is finite because the exponential dominates the polynomial factor. Thus the exponential measures of \cref{thm:exponential-damping} are continuous under capacity convergence. Their tails satisfy the useful bound
\[
 \nu_{p,u}(\{u<-T\})
 \leq\mathrm{e}^{-pT/2}\nu_{p/2,u}(X)
 \leq C_p\mathrm{e}^{-pT/2}.
\]
Polynomial weights give a more explicit form of the general estimate.


\section{Uniform tails and undamped convergence}

Removing the weight requires uniform control of the mass on deep sublevels of the approximating potentials. The following condition supplies both a finite limit measure and weak convergence. It can equivalently be expressed by an increasing integrability weight.

\begin{theorem}\label{thm:undamped-tail-continuity}
Let $u_j,u\in\PSH(X,\omega)$ be nonpositive and suppose $u_j\to u$ in capacity. Assume
\[
 \lim_{T\to\infty}\sup_j\mu_{u_j}(\{u_j<-T\})=0.
\]
Then $\mu_u$ is finite and $\mu_{u_j}\rightharpoonup\mu_u$.
\end{theorem}
\begin{proof}
Choose $T_0$ so that every tail beyond $T_0$ has mass at most one. By \cref{lem:sublevel-mass-growth}, $\mu_{u_j}(X)\leq C(1+T_0)^{r_n}+1$. Thus the tail condition already gives a uniform bound for the total masses.

Choose a continuous nonincreasing function $\eta:[0,+\infty)\to[0,1]$, equal to one on $[0,1]$ and zero on $[2,+\infty)$. Put $\eta_T(s)=\eta(s/T)$. The fixed-truncation argument in \cref{thm:damped-continuity}, now with coefficient $\eta_T(-u_j)$, gives
\[
 \eta_T(-u_j)\mu_{u_j}\rightharpoonup\eta_T(-u)\mu_u
 \quad\text{for every fixed }T>0.
\]
The integrals on the left are bounded by this uniform total mass. Since $\eta_T(-u)$ increases to one off the pole set as $T\to\infty$, monotone convergence proves $\mu_u(X)<+\infty$.

For a continuous test function $f:X\to\R$, the error made by inserting $\eta_T(-u_j)$ is at most
$\|f\|_\infty\mu_{u_j}(\{u_j<-T\})$, uniformly small for large $T$. For the limit potential the corresponding error tends to zero because $\mu_u$ is finite and its pole set has zero mass. Pass first $j\to\infty$ with $T$ fixed, and then $T\to\infty$.
\end{proof}

\Needspace{8\baselineskip}
The tail condition has an equivalent formulation using an increasing weight. Let $\chi:[0,+\infty)\to[1,+\infty)$ be continuous, increasing and unbounded. For a nonpositive potential define
\[
 J_{\chi,\omega}(u)=\int_X\chi(-u)\,\mathrm{d}\mu_u,
\]
with the value $+\infty$ allowed. A uniform bound for this functional forces the sublevel tails to vanish because $\chi(T)\to+\infty$.

\begin{corollary}\label{cor:integrability-continuity}
If $u_j\to u$ in capacity and $\sup_jJ_{\chi,\omega}(u_j)\leq C$, then
\[
 J_{\chi,\omega}(u)\leq C,\qquad
 \mu_{u_j}\rightharpoonup\mu_u.
\]
\end{corollary}
\begin{proof}
Since $\chi\geq1$, the total masses are at most $C$, and
\[
 \mu_{u_j}(\{u_j<-T\})\leq C/\chi(T).
\]
Thus \cref{thm:undamped-tail-continuity} applies. For the integrability assertion, use its fixed-truncation argument with the bounded continuous coefficient
$\chi(-u_j)\eta_T(-u_j)$. Its integral converges at each fixed $T$ and is at most $C$. Let $T\to\infty$ by monotone convergence.
\end{proof}

\begin{proposition}\label{prop:tail-integrability-equivalence}
For a family of nonpositive potentials, the uniform tail condition in \cref{thm:undamped-tail-continuity} is equivalent to a uniform bound on $J_{\chi,\omega}$ for some continuous increasing unbounded $\chi\geq1$.
\end{proposition}
\begin{proof}
One implication was proved in the preceding corollary. Conversely, write the family as $(u_j)_{j\in J}$, where $J$ is its index set. The tail condition and \cref{lem:sublevel-mass-growth} give a uniform bound for the total masses, as in \cref{thm:undamped-tail-continuity}. Choose numbers $T_k\to+\infty$, with $T_{k+1}>T_k+1$, such that the mass of $\{-u_j>T_k\}$ is at most $2^{-k}$ for every member of the family. Let $r_k$ be a continuous nondecreasing function, zero on $[0,T_k]$ and one on $[T_k+1,+\infty)$. Define
\[
 \chi(s)=1+\sum_{k\geq1}r_k(s).
\]
The sum is locally finite, so $\chi$ is continuous, increasing and unbounded. Tonelli's theorem gives
\[
 J_{\chi,\omega}(u_j)
 \leq\mu_{u_j}(X)+\sum_{k\geq1}\mu_{u_j}(\{-u_j>T_k\})
 \leq\sup_j\mu_{u_j}(X)+1.
\]
\end{proof}

The following example exhibits mass concentrating at a pole of the limiting potential.

\Needspace{14\baselineskip}
\begin{example}\label{ex:p1-raw-tail}
The tail condition cannot be omitted, even on a compact K\"ahler curve. On $X=\mathbb P^1$, let $z$ be the standard affine coordinate and let $\delta_0$ denote the unit Dirac measure at its origin. Take the Fubini--Study form
\[
 \omega=\ddc\log(1+|z|^2),\qquad \int_X\omega=1,
\]
and the global nonpositive potential
\[
 u(z)=\log\frac{|z|^2}{1+|z|^2},\qquad
 u_j=\max(u,-j).
\]
At infinity, in the coordinate $\zeta=1/z$, the function $u$ is $-\log(1+|\zeta|^2)$, so it is smooth there. The Poincar\'e--Lelong formula gives
\[
 \omega+\ddc u=\delta_0,
 \qquad \mathrm{MA}^{\mathrm{ext}}_\omega(u)=0.
\]
The second identity follows because the extended measure discards the pole set and the current vanishes on its complement.

Each $u_j$ is bounded and $\omega$-plurisubharmonic. Its measure
$\mu_j=\omega+\ddc u_j$ has mass one by Stokes' theorem. On the open set $\{u>-j\}$, locality gives $\mu_j=0$. Hence
\[
 \operatorname{supp}\mu_j\subset\{u\leq-j\}
 =\left\{|z|^2\leq\frac{\mathrm{e}^{-j}}{1-\mathrm{e}^{-j}}\right\}.
\]
These discs shrink to the origin, so $\mu_j\rightharpoonup\delta_0$.
\Needspace{5\baselineskip}
On their supports, $u_j=-j$. Consequently, for $j>T$,
\[
 \mu_j(\{u_j<-T\})=1,
 \qquad \mathrm{e}^{pu_j}\mu_j=\mathrm{e}^{-pj}\mu_j\rightharpoonup0
 \quad(p>0).
\]
Nevertheless, $u_j\to u$ in capacity. Indeed, outside any neighborhood of the origin the functions eventually agree, and neighborhoods of this capacity-zero point can be chosen with arbitrarily small capacity. Thus bounded mass and capacity convergence do not prevent the undamped measures from acquiring mass on a pole of the limit. The damped limit agrees with the extended measure, as predicted by \cref{thm:damped-continuity}.
\end{example}

The preceding example has a singular limit. When bounded mass fails, the smooth sequence of \cref{cor:smooth-prescribed-masses} has a smooth limit and still loses continuity of its total mass. The loss occurs entirely at deep sublevels. For a finite signed measure $\tau$, write $\|\tau\|_{\mathrm{TV}}$ for its whole total variation norm.

\Needspace{16\baselineskip}
\begin{corollary}\label{cor:smooth-capacity-mass-defect}
Assume $v_+(\omega)=+\infty$ and put $V_\omega=\int_X\omega^n$. Fix a real number $b>V_\omega$ and a sequence $0<\varepsilon_j<1$ tending to zero. The smooth potentials $w_j$ in \cref{cor:smooth-prescribed-masses}, chosen with mass $b$, converge to zero in capacity. For every $T>0$,
\[
 \left\|\mathbf1_{\{w_j\geq-T\}}(\omega+\ddc w_j)^n-\omega^n
 \right\|_{\mathrm{TV}}\longrightarrow0,
 \qquad
 \int_{\{w_j<-T\}}(\omega+\ddc w_j)^n\longrightarrow b-V_\omega.
\]
After passage to a subsequence, there is a positive measure $\sigma$ of mass $b-V_\omega$ such that
\[
 (\omega+\ddc w_j)^n\rightharpoonup\omega^n+\sigma,
 \qquad
 \mathbf1_{\{w_j<-T\}}(\omega+\ddc w_j)^n\rightharpoonup\sigma
 \quad\text{for every }T>0.
\]
\end{corollary}
\begin{proof}
The construction in \cref{cor:smooth-prescribed-masses} gives $w_j=s_j u_j$, where $u_j$ is smooth and $\omega$-plurisubharmonic, $\sup_Xu_j=0$ and $0<s_j<\varepsilon_j$. The normalized sublevel estimate \eqref{eq:global-capacity-tail} gives, for every fixed $\delta>0$ and all sufficiently large $j$,
\[
 \operatorname{Cap}_\omega(\{|w_j|>\delta\})
 =\operatorname{Cap}_\omega(\{u_j<-\delta/s_j\})
 \leq Cs_j/\delta\longrightarrow0.
\]
Put $\mu_j=(\omega+\ddc w_j)^n$. Fix $T>0$ and write $E_{j,T}=\{w_j\geq-T\}$. The lower bound in \cref{cor:smooth-prescribed-masses} gives
\[
 \mathbf1_{E_{j,T}}\mu_j\geq
 \lambda_{j,T}:=(1-\varepsilon_j)^n\mathbf1_{E_{j,T}}\omega^n.
\]
Since $w_j\to0$ in $L^1$, the smooth volume of $X\setminus E_{j,T}$ tends to zero. Hence $\lambda_{j,T}\to\omega^n$ in total variation. For every $p>0$, \cref{thm:damped-continuity} gives $\int_X\mathrm{e}^{pw_j}\,\mathrm{d}\mu_j\to V_\omega$, and therefore
\[
 \limsup_j\mu_j(E_{j,T})
 \leq\mathrm{e}^{pT}\lim_j\int_X\mathrm{e}^{pw_j}\,\mathrm{d}\mu_j
 =\mathrm{e}^{pT}V_\omega.
\]
Let $p\downarrow0$. Together with the lower bound this proves $\mu_j(E_{j,T})\to V_\omega$. Positivity now yields
\[
 \begin{split}
 \|\mathbf1_{E_{j,T}}\mu_j-\omega^n\|_{\mathrm{TV}}
 &\leq\mu_j(E_{j,T})-\lambda_{j,T}(X)
       +\|\lambda_{j,T}-\omega^n\|_{\mathrm{TV}}\\
 &\longrightarrow0.
 \end{split}
\]
Subtracting this shallow mass from $\mu_j(X)=b$ proves the tail identity.

The positive measures $\mu_j$ have fixed mass on compact $X$, so a subsequence converges weakly to a measure $\nu$ of mass $b$. The inequality $\mu_j\geq(1-\varepsilon_j)^n\omega^n$ gives $\nu\geq\omega^n$. Set $\sigma=\nu-\omega^n$. For each $T>0$, subtract the already established total-variation limit of the shallow restrictions from $\mu_j\rightharpoonup\nu$. This gives the last assertion along the same subsequence for every $T$.
\end{proof}

The currents in \cref{cor:smooth-capacity-mass-defect} also converge in coefficient $L^1$, and their Gauduchon traces converge in total variation. The argument identifies the mass of $\sigma$, but does not determine its support or whether it charges pluripolar sets.

\section{Damping for a smooth background}

For a general smooth background $\theta$, constant truncation may fail to preserve $\theta$-plurisubharmonicity. A larger Hermitian form makes the truncations positive and dominates the measure; the local expansion of the $\theta$-product then gives continuity at each fixed truncation level.

\begin{theorem}\label{thm:general-background-damping}
Let $\theta$ be any smooth real $(1,1)$-form. For nonpositive $u\in\PSH(X,\theta)$ put $\mu_{\theta,u}=\mathrm{MA}^{\mathrm{ext}}_\theta(u)$. The finite damping, capacity estimates and capacity-continuity statements of this chapter, including \cref{prop:general-damping-weights}, hold for these measures, with constants depending on $\theta$ and a fixed Hermitian metric. In particular, for $p>0$ and $q>r_n$,
\[
 \int_X\mathrm{e}^{pu}\,\mathrm{d}\mu_{\theta,u}
 \leq C(1+p^{-1})^{r_n},\qquad
 \int_X(1-u)^{-q}\,\mathrm{d}\mu_{\theta,u}\leq C_q.
\]
The undamped tail criterion and the integrability criterion also hold with $\mu_{\theta,u}$ in place of $\mu_u$.
\end{theorem}
\begin{proof}
Choose $C_0\geq1$ such that $\theta\leq\Omega:=C_0\omega$. Then $u\in\PSH(X,\Omega)$ and \cref{thm:arbitrary-extended-ma} gives
\[
 0\leq\mu_{\theta,u}\leq\mathrm{MA}^{\mathrm{ext}}_\Omega(u).
\]
All the mass, tail and capacity bounds follow from the already proved estimates for $\Omega$. To express them in the original capacity, note that for every Borel set $E\subset X$,
\[
 \operatorname{Cap}_\omega(E)
 \leq\operatorname{Cap}_\Omega(E)
 \leq C_0^n\operatorname{Cap}_\omega(E).
\]
The first inequality follows by inclusion of the test potentials and positivity of the added background. For the second, if $0\leq v\leq1$ is $\Omega$-plurisubharmonic, then $v/C_0$ is an admissible $\omega$-test potential and
$(\Omega+\ddc v)^n=C_0^n(\omega+\ddc(v/C_0))^n$.

For continuity, let nonpositive $u_j,u\in\PSH(X,\theta)$ converge in capacity. Use the fixed-truncation identity in \cref{thm:damped-continuity}. Its truncations are bounded and $\Omega$-plurisubharmonic. Their signed $\theta$-products are continuous under capacity convergence by the local expansion in \cref{prop:bounded-products}, also with the bounded quasi-continuous coefficients used above. The cutoff vanishes at the truncation level, giving exactly the restricted positive measure. The uniform tails supplied by domination remove the cutoff. Domination also supplies the sublevel estimate, so the preceding proofs of the undamped tail and integrability criteria apply verbatim.
\end{proof}
